El principal objetivo del proyecto es el desarrollo de una herramienta que identifica y separa los diferentes elementos en una obra musical.

Para facilitar el llevar a cabo el objetivo principal, facilita la tarea dividirlo en batallas más pequeñas. Se pueden resumir en dos grupos a los que se suma un objetivo secundario:

\begin{itemize}
    \item \textbf{1. Objetivo inicial:} separación de dos fuentes reduciendo la tarea a la identificación de únicamente vocales y acompañamiento.
    
    Este objetivo se daría por logrado al conseguir una implementación que diese como salida dos archivos de audio donde se discerniesen dos elementos diferentes provenientes de la mezcla original, y que, en conjunto, sumasen esta. Este objetivo se daria por cumplido al tener un sistema implementado cuya salida consistiese en las dos fuentes de una pista con dos elementos con una distinción clara, sin presencia de interferencias, sin presencia de ruido y que en conjunto sumasen la pista original.
    \item \textbf{2. Objetivo a largo plazo:} una vez cumplido el objetivo inicial, el siguiente paso es llegar a la separación de cuatro fuentes recogiendo vocales , bajo , baterías y el resto del acompañamiento identificado como "\textit{resto}".
    
    Este objetivo se daría por cumplido al conseguir una implementación en código que diese como salida cuatro archivos de audio comprendiendo vocales,bajo,baterías y resto de elementos, con una distinción clara de los elementos, sin presencia de interferencias, sin presencia de ruido y que en conjunto sumasen la pista original.
\end{itemize}

\textbf{Objetivo secundario:} Adicionalmente,y, aunque no estaba dentro de los objetivos iniciales y se tomó mas como una propuesta a realizar en caso de que lo permitiese la planificación, se planteó desarrollar una interfaz gráfica muy sencilla que permite cargar audios de prueba y obtener sus \textit{stems}\footnote{\textit{STEMS} es el término anglosajón que refiere a las sub-pistas que componen una mezcla musical ( en castellano, fuentes)} separados, así como almacenar los resultados en forma de archivos de audio y gráficas con cierto enfoque didáctico.


En resumen, podríamos agrupar los objetivos en:

\begin{itemize}
    \item \textbf{1. Objetivo principal: } llegar a un estado del proyecto donde se genere una salida que presente una diferencia respecto a la pista de entrada y que identifique al menos superficialmente el elemento objetivo.
    \item \textbf{2. Objetivos técnicos: } implementar un proceso funcional de entrenamiento para las diferentes funciones. Evaluar los resultados y proporcionar una comparativa de las funciones.
    \item \textbf{3. Objetivo secundario de la interfaz: } facilitar el uso, dar comprensión para principiantes en el campo y dar visualización del sistema.
\end{itemize}
